\chapter{Grading: From Weight Maps to the Where-Clause}
\label{ch:weight}
\label{ch:grading}

\section{What is being built, and in what order}

A rewrite rule's left-hand side is a type. Terms inhabit it; the rule fires at its
inhabitants. This chapter is about what happens when one attaches a number to the
firing, and it builds the attachment twice.

The first construction is the \emph{weight map}. Given a continued interactive GSLT $G$,
the weight endofunctor $\Wgt_{\Semi}$ produces a theory in which every base rule carries a
finite table
\[
  \varphi_1 \mapsto v_1, \quad \ldots, \quad \varphi_m \mapsto v_m
\]
from formulae refining its left-hand side to values in a semiring $\Semi$. A term's
redexes are classified by which refinement they satisfy, and the values turn the set of
enabled redexes into a distribution. Where $\Cost$ adjoins a stack of tokens and $\Hist$
adjoins a word of events, $\Wgt_{\Semi}$ adjoins a table of rates --- and, like both of
them, it adjoins it \emph{to the state}. Sections~\ref{wt:sec:partition}
to~\ref{wt:sec:spim} build this and prove what it gives.

The second construction is the \emph{graded where-clause}, and it is the one the rest of
the book uses. It is stated in a sentence: \emph{the \texttt{where} clause of the
for-comprehension is valued in a commutative semiring and is evaluated once per candidate
match.} There is no table, no key set and no partition. A programmer writes the weight
where the program branches, in the same language the program is written in, and the
value may depend on what was matched.
Sections~\ref{gw:sec:choice} to~\ref{gw:sec:contention} build this.

The two are not rivals, and the order is deliberate. The weight map is the easier object
to see: it is a table, the table is attached to a rule, and a simulator reads it. Having
seen it, the reader can watch it dissolve. Theorem~\ref{gw:thm:dnf} says that a weight map
is exactly a graded clause of one particular shape --- a sum over a partition --- so that
everything Sections~\ref{wt:sec:partition}--\ref{wt:sec:spim} prove about weight maps is
a theorem about the \emph{linear} fragment of the clause language, and the discipline those
sections had to impose is the price of staying in that fragment. The reader who already
knows the chapter's destination may read the first half as an extended example. The reader
who does not will find that the second half has been argued for before it is stated.

Two things distinguish both constructions from the literature on stochastic process
calculi, and both are consequences of taking the weights from a logic the theory
generates rather than from a vocabulary the modeler invents.

\paragraph{Weights depend on formulae, not names.} In Priami's stochastic
$\pi$-calculus~\cite{priami} and in Phillips and Cardelli's Stochastic Pi Machine
(SPiM)~\cite{spim}, a rate is attached to a \emph{channel}. That is one refinement among
many: ``this redex is a communication on $n$'' is a particular formula, and there is no
reason a modeler should be confined to it. Once weights are given by formulae, a rate may
depend on the shape of the payload, on the \emph{namespace} a channel belongs to --- which
in a reflective calculus is real structure, reachable through the name predicate --- or
on any combination of these. Channel identity is recovered as the degenerate case, and
Theorem~\ref{wt:thm:spim} says exactly how degenerate it is.

\paragraph{Weights are state, not declaration.} Rates change during execution. SPiM does
not permit this, and the restriction is not incidental to its design but constitutive of
it, since a static rate function is what makes the underlying object a chain over terms
alone. The weight map recovers Markovianity by putting the table \emph{in} the state, so
the chain runs over configurations $(P,\dec)$ rather than over terms
(Theorem~\ref{wt:thm:markov}). The where-clause recovers it more simply: a clause is
re-evaluated against the current term at every candidate, so whatever a rate depends on
is already in the term, and the term is the state
(Proposition~\ref{gw:prop:plastic}). Chapter~\ref{ch:spiking} depends on one or the other,
since plasticity is nothing other than a rate that changes.

\begin{remark}[Where the formulae come from]
\label{wt:rmk:whence}
The formulae used, as keys in the first construction and as the crisp sort of the
clause language in the second, are those of the logic the OSLF construction generates
from the theory's own presentation: one structural connective per term former of
$\Sigma$, one context-labeled modality $\langle K \rangle$ per rewrite rule together with
a choice of redex position, and a propositional layer stratified by the categorical
strength the target specification supplies. That construction is
Chapter~\ref{ch:hypercube}, and the transition system the modalities range over is
Chapter~\ref{sec:lts-hml}. Nothing in the present chapter needs the details; it needs
only that the vocabulary of refinements is \emph{generated with the theory} rather than
designed alongside it, so that a modeler who writes down a language has, at no further
cost, the language in which to price it. The one place the details do intrude is
Observation~\ref{wt:obs:closure}, which is about what a weak target specification cannot
say --- and Corollary~\ref{gw:cor:nopartition}, which is about why, in the end, it does
not have to.
\end{remark}

\section{Keys, and the partition discipline}
\label{wt:sec:partition}

Fix a base rule $r : L \rsq R$. Write $L^{\sharp}$ for the structural formula
characterizing the left-hand side pattern --- for the rho communication rule, ``an input
on some $n$ in parallel with an output on the same $n$'' with $n$ existentially bound.

\begin{definition}[Refinement]
\label{wt:def:refinement}
A formula $\varphi$ \emph{refines} $L^{\sharp}$ when $\sat \varphi \Rightarrow L^{\sharp}$.
The refinements of $L^{\sharp}$, ordered by entailment, form the \emph{refinement lattice}
of the rule.
\end{definition}

A key set is a finite subset of that lattice. Three requirements are imposed on it, and
the first is the substantial one.

\begin{definition}[Partition discipline]
\label{wt:def:partition}
A key set $\Phi_r = \{\varphi_1, \ldots, \varphi_m\}$ is a \emph{partition of $r$} when
\begin{enumerate}[label=(P\arabic*),leftmargin=3em]
  \item \emph{Exclusivity}: $\sat \neg(\varphi_i \wedge \varphi_j)$ for $i \neq j$;
  \item \emph{Exhaustiveness}: $\sat L^{\sharp} \Rightarrow \bigvee_i \varphi_i$.
\end{enumerate}
\end{definition}

\begin{definition}[Admissible key set]
\label{wt:def:reqs}
A key set is \emph{admissible} when it is a partition and in addition
\begin{enumerate}[label=(R\arabic*),leftmargin=3em]
  \item\label{wt:req:dec} \emph{Decidability}: each $\varphi_i$ lies in a fragment for
        which model checking terminates on the terms of interest;
  \item\label{wt:req:loc} \emph{Locality}: the modal depth of each $\varphi_i$ is bounded
        by some $\alpha_r < \infty$ fixed with the rule;
  \item\label{wt:req:str} \emph{Structurality}: each $\varphi_i$ lies in the structural
        fragment, carrying no behavioral modality $\langle K \rangle$, no greatest fixed
        point, and no hidden-name quantifier.
\end{enumerate}
\end{definition}

\ref{wt:req:dec} is forced by execution: a simulator must classify every enabled redex
before it can sample, so a key whose evaluation may diverge is a key that stops the run.
\ref{wt:req:loc} and \ref{wt:req:str} are forced by \emph{efficient} execution, and they
are restrictions on \emph{keys} alone --- the modalities and the fixed point remain
available for stating and checking properties of the resulting chain. What they may not
do is determine a rate. A rate depending on what a redex can do \emph{next} cannot be
recomputed locally when something elsewhere in the term changes, so an implementation
admitting modal keys must reclassify globally at every step. The clause language of
Section~\ref{gw:sec:clause} will relax exactly this, at exactly this price.

\begin{proposition}[Classification]
\label{wt:prop:classify}
Let $\Phi_r$ be a partition of $r$ and let $\Red_r(P)$ be the set of redexes of $r$ in
$P$, each a pair $(k,\sigma)$ of a position and a matching substitution. Then
\[
  \class_r : \Red_r(P) \longrightarrow \Phi_r,
  \qquad
  \class_r(k,\sigma) = \text{the unique } \varphi_i \text{ with } L\sigma \sat \varphi_i
\]
is a total function, for every term $P$.
\end{proposition}

\begin{proof}
$(k,\sigma) \in \Red_r(P)$ gives $P \equiv k[L\sigma]$, so $L\sigma \sat L^{\sharp}$ by
construction of $L^{\sharp}$. Exhaustiveness gives some $i$, which is totality;
exclusivity gives at most one, which is single-valuedness. Well-definedness on
$\equiv$-classes holds because satisfaction is $\equiv$-invariant, the generated
connectives being generated from the signature modulo the equations.
\end{proof}

Proposition~\ref{wt:prop:classify} is the whole reason for the discipline, and what it
buys is visible in what fails without it. Without exhaustiveness the weighting is
\emph{partial}: some enabled redexes have no rate, and one must invent a convention ---
fire them at rate zero, at a default rate, or refuse to run --- that is not part of the
specification. Without exclusivity the weighting is \emph{multi-valued}, and one must
invent an aggregation: take the most specific key if a unique minimum exists, sum the
values of all satisfied keys, take the maximum. Each convention is defensible, and each
yields a different simulator from the same specification, which is precisely the situation
a specification exists to prevent. Hold on to the second of those conventions --- sum the
values of all satisfied keys --- because Section~\ref{gw:sec:dnf} will show that it is not
a convention at all but the semantics, once one stops asking for a table.

\begin{remark}[Completing a partition is cheap]
\label{wt:rmk:default}
The discipline is less onerous than it looks. Given any pairwise-exclusive family
$\{\varphi_1,\ldots,\varphi_m\}$, adjoining
$\varphi_{\bot} := L^{\sharp} \wedge \neg\bigvee_i \varphi_i$ yields a partition, and a
modeler wanting the unlisted redexes inert sets its value to $0$. Exclusivity is the
real constraint; exhaustiveness is a completion. In practice, keys of the form ``the
redex is a communication on channel $n$'', indexed by distinct $n$, are exclusive for
free.
\end{remark}

\subsection{What stating the partition costs}
\label{wt:sec:closure}

Definition~\ref{wt:def:partition} writes (P1) and (P2) as entailments. If they are to be
stated \emph{in the generated logic itself}, that logic must carry negation, conjunction
and disjunction, and a target specification supplying only the finite-limits fragment
will not do --- finite limits supply $\top$ and $\wedge$ and no more
(Remark~\ref{ob:rem:stratification}).

This is a hypothesis and not a theorem about weighting as such. A key family drawn from a
positive fragment may perfectly well \emph{be} pairwise disjoint and jointly exhaustive;
what a weak logic prevents is not the holding of (P1) and (P2) but their \emph{expression}
as entailments between formulae of the language the keys are written in. The alternative
is to discharge them in the metatheory --- by an elaborator checking them on the instances
that arise, or by a proof about a schema of keys --- and the weighting is then perfectly
well defined. Both branches cost something, and the second cost is the interesting one.

\begin{observation}[The cost of an external partition]
\label{wt:obs:closure}
If (P1) and (P2) are discharged in the metatheory rather than in the generated logic, the
specification is no longer \emph{closed}: whether a key set is admissible becomes a
judgment the object language cannot express. Two implementations conforming to the same
written specification may then disagree about admissibility, hence --- by
Proposition~\ref{wt:prop:classify} --- about the classification, hence about the
distribution sampled. That is exactly the failure the partition discipline was introduced
to prevent, relocated one level up rather than removed.
\end{observation}

So a weighted theory built on weight maps must choose between a Boolean target
specification and an external admissibility judgment, and it should say which. This is a
small instance of a pattern the Prestige will make central: a construction that looks
self-contained turns out to have a seam, and the seam is where something the model cannot
state has to be supplied from outside (Chapter~\ref{ch:notation}). It is worth noticing
that the seam here is \emph{visible} --- one can point at the entailment that has to be
checked somewhere else --- and that this is the good case. It is also worth noticing, in
advance, that this particular seam is an artifact of asking for a table, and that
Corollary~\ref{gw:cor:nopartition} closes it.

\section{The endofunctor}
\label{wt:sec:functor}

\begin{definition}[Weight map]
\label{wt:def:wmap}
Let $G$ be a continued interactive GSLT with an admissible key set $\Phi_r$ for each base
rule $r$. A \emph{weight map} is a function
\[
  \dec : \textstyle\coprod_r \Phi_r \longrightarrow \Semi .
\]
Write $\Wmaps$ for the set of weight maps. Each $\Phi_r$ is finite and there are finitely
many rules, so $\Wmaps \cong \Semi^N$ for a fixed $N$.
\end{definition}

\begin{definition}[Configuration]
\label{wt:def:config}
A \emph{configuration} is a pair $\mathfrak{c} = (P, \dec)$ with $P$ a term modulo
$\equiv$ and $\dec \in \Wmaps$. Write $\Cfg$ for the set of configurations.
\end{definition}

The phrase ``a map instance is part of a state specification'' is
Definition~\ref{wt:def:config}, and everything in this half of the chapter turns on it.

\begin{definition}[Augmented rules]
\label{wt:def:augmented}
An \emph{augmented base rule} is a base rule $r : L \rsq R$ together with, for each
$\varphi_i \in \Phi_r$, an initial value $v_i \in \Semi$ and an \emph{update function}
$u_i : \Wmaps \to \Wmaps$ on the whole map. An \emph{augmented context rule} is a context
rule together with a \emph{geometric factor} $\gfac_c \in \Rnn$ and a fold function
$f_c$ combining the maps arriving from its sibling subterms. The initial weight map is
$\dec_0(\varphi_i) = v_i$.
\end{definition}

\begin{construction}[The weight endofunctor]
\label{wt:con:functor}
$\Wgt_{\Semi}$ sends a continued interactive GSLT $G$, equipped with admissible key sets,
to the theory $\Wgt_{\Semi}(G)$ whose configurations are the pairs of
Definition~\ref{wt:def:config} and whose rules are the augmented rules of
Definition~\ref{wt:def:augmented}: a step is
\[
  (P, \dec) \rewrite (P', \dec'),
  \qquad
  P' \equiv k[R\sigma],
  \qquad
  \dec' = f_{c_1}\bigl(\cdots f_{c_p}(u_i(\dec))\cdots\bigr),
\]
for $(k,\sigma)$ a redex of $r$ with $\class_r(k,\sigma) = \varphi_i$ and $k$ composed of
the context rules $c_1, \ldots, c_p$ from the root inwards, the folds applied from the
redex outwards. In the common case where every $f_c$ projects onto its second argument ---
``the context does not modify the map'' --- this is $\dec' = u_i(\dec)$, and we write it so
below. On morphisms, $\Wgt_{\Semi}$ acts as the identity on the underlying map of theories
and by reindexing on the key sets.
\end{construction}

The parallel with the two preceding chapters is exact, and it is the reason this chapter
sits where it does. Each of the three constructions adjoins a component to the state and
makes the rules act on it: a stack that drains, a word that accrues, a table that is
rewritten. Each leaves the underlying interaction structure alone. And each is a genuine
addition rather than a closure operation --- the table after a step is not in general the
table before it.

\begin{remark}[Two adjunctions and one open question]
\label{wt:rmk:notamonad}
$\Cost$ and $\Hist$ are monads, resolved by free--forgetful adjunctions
(Proposition~\ref{ob:prop:cost-adjunction}, Proposition~\ref{ob:prop:hist-adjunction}).
We claim no such thing for $\Wgt_{\Semi}$. There is an evident candidate unit: the map
sending every key to $1_{\Semi}$, which embeds $G$ as the uniformly weighted theory in
which every enabled redex is equally likely, and which is unital in the right way for the
same reason the empty stack and the empty word are. Since a unit is available it might
well be interesting to ask whether a multiplication exists --- what it would mean to
flatten a weighting of a weighting, and whether the flattening is canonical rather than
chosen. We do not ask it here. Nothing below requires more than an endofunctor, and
stating the construction is the point.
\end{remark}

\begin{remark}[Why context rules carry a factor and not a propensity]
\label{wt:rmk:geometric}
The obvious alternative to Definition~\ref{wt:def:augmented} gives every rule, base and
context alike, a value, and lets the simulator select among all of them. It
double-counts. A step is a base rule firing \emph{somewhere}; the context rules composing
the path from the root to that somewhere are the \emph{address} of the firing, not
additional events. Structural rules are addressing, not dynamics, and incur no primitive
cost. So the context rules contribute multiplicatively, by
$\gfac(k) = \prod_q \gfac_{c_q}$ over the composition, and setting every $\gfac_c = 1$
recovers ``addressing is free'' exactly, which is the default.

The reason to permit $\gfac_c \neq 1$ is that it is the natural home of a \emph{spatial}
bias. $\gfac$ weights an edge of the reduction graph by \emph{where} in the term the
interaction surfaces are; $\dec(\varphi)$ weights it by \emph{what} is being transferred.
That is a geometry/matter split appearing here as an implementation detail, and
Chapter~\ref{sec:discussion} says what would have to be true for it to be more than one.
\end{remark}

\section{Propensity}
\label{wt:sec:propensity}

The propensity of a rule at a refinement is where the construction is most easily got
wrong, because four independent factors are involved and the literature usually presents
some of them fused.

\begin{definition}[Propensity]
\label{wt:def:propensity}
Let $\mathfrak{c} = (P, \dec)$, let $r$ be a base rule and $\varphi \in \Phi_r$. Put
$K(r,\varphi,P) = \{(k,\sigma) \in \Red_r(P) : \class_r(k,\sigma) = \varphi\}$. The
\emph{propensity} of $(r,\varphi)$ at $\mathfrak{c}$ is
\[
  \prop(r,\varphi,\mathfrak{c})
  \;=\;
  \dec(\varphi) \cdot \!\!\sum_{(k,\sigma) \in K(r,\varphi,P)}\!\! \gfac(k) \cdot \chi(r,k,\sigma),
\]
where $\chi \in \{0,1\}$ is the funding gate of Definition~\ref{wt:def:gate}, identically
$1$ in the absence of cost accounting. The \emph{total propensity} is
$\propz(\mathfrak{c}) = \sum_r \sum_{\varphi \in \Phi_r} \prop(r,\varphi,\mathfrak{c})$.
\end{definition}

The four factors, and what each answers:

\begin{description}[leftmargin=1.6em,style=nextline]
  \item[$\dec(\varphi)$ --- the rate constant.] What kind of interaction this is, priced.
        Dynamic: this is the component the update functions rewrite.
  \item[$\gfac(k)$ --- the geometric factor.] Where the interaction is, priced. Default
        $1$.
  \item[The cardinality of the sum --- multiplicity.] How many ways this interaction is
        available. This is Gillespie's combinatorial count of distinct reactant
        tuples~\cite{gillespie1977exact}, and it is the factor most often omitted in
        implementations, with the consequence that a term with fifty available redexes of
        a kind fires them no faster than a term with one.
  \item[$\chi$ --- the funding gate.] Whether the interaction can be afforded.
\end{description}

Multiplicity deserves a word, because it is the one place where the associative--commutative
structure of the interaction constructor does real work. Redexes are counted as
\emph{selections} from the bag, not identified up to $\equiv$. Two indistinguishable
messages on the same channel offer two ways for a receipt to fire, and a construction that
identified them would run the system at half rate. This is not a subtlety of the
implementation; it is the difference between a truncated Poisson stationary distribution
and a geometric one. The clause language will make the same point in its own terms:
candidates are indexed by occurrences, not by terms (Definition~\ref{gw:def:cand}).

\begin{definition}[Funding gate]
\label{wt:def:gate}
For a redex $(k,\sigma)$ of $r$, let $\Delta(r,k,\sigma)$ be the demand it incurs and
$\Sigma(P,k)$ the supply available at that location, in the sense of
Chapter~\ref{ch:costmonad}. Put $\chi(r,k,\sigma) = 1$ when $\Sigma(P,k)$ funds
$\Delta(r,k,\sigma)$ and $0$ otherwise.
\end{definition}

\begin{proposition}[Cost gates the support, weight distributes over it]
\label{wt:prop:gate}
Unfunded redexes have propensity zero and contribute nothing to $\propz$, while among
funded redexes the distribution is unchanged from the cost-free case up to normalization.
\end{proposition}

\begin{proof}
$\chi$ appears as a multiplicative factor taking values in $\{0,1\}$; a zero factor
removes the term from the sum, and the remaining terms are untouched.
\end{proof}

Three consequences are worth recording, and they are the reason this chapter belongs
beside Chapter~\ref{ch:costmonad} rather than anywhere else in the book.

\emph{The gate is decidable, so the simulator stays exact.} Were funding a real-valued
penalty rather than a Boolean, the propensity would depend on an optimization and the
process would no longer be exactly simulable. The decidability of the funding judgment is
doing real work here.

\emph{Exhaustion is a halting condition, not a deadlock.} $\propz = 0$ when every enabled
redex is unfunded. The chain has reached an absorbing state, which is the correct reading:
a system that cannot pay for any of its available transitions has stopped, and the waiting
time to the next event is infinite rather than undefined.

\emph{The two decorations answer different questions.} Cost decides whether a step may
happen; weight decides how fast and how often. Separating them is what lets a model say
that a strong synapse on a starving neuron is silent --- a statement neither decoration
makes alone, and one Chapter~\ref{ch:spiking} needs.

\section{The simulator is a functor out of the image}
\label{wt:sec:simulator}

$\Wgt_{\Semi}(G)$ is still a rewrite theory: it has terms, equations and rules, and it
steps. Turning it into something one can sample from is a further operation, and the two
should not be run together.

\begin{definition}[The weighted coalgebra]
\label{wt:def:coalgebra}
Let $\Semi^{(-)}$ denote the free-$\Semi$-semimodule construction, $\Semi^{(X)}$ being the
finitely supported functions $X \to \Semi$. Evaluating
Definition~\ref{wt:def:propensity} pointwise gives
\[
  \gamma : \Cfg \longrightarrow \Semi^{(\Cfg)},
  \qquad
  \gamma(\mathfrak{c})(\mathfrak{c}')
  = \!\!\sum_{\substack{(r,k,\sigma) \\ \text{stepping } \mathfrak{c} \to \mathfrak{c}'}}\!\!
    \dec(\class_r(k,\sigma)) \cdot \gfac(k) \cdot \chi(r,k,\sigma).
\]
\end{definition}

The sum runs over \emph{derivations} reaching $\mathfrak{c}'$, not merely over the target.
Assembling the weightings of a fixed theory into a category and taking the Grothendieck
construction gives the fibered form of this assignment, and the simulator is a functor out
of it. We record that to fix the picture and do not need it below.

\begin{remark}[Simulator and interpreter are different objects]
\label{wt:rmk:simulator}
It is worth saying plainly that $\gamma$ is not what a running system does.
An interpreter advances one actual state; a simulator explores how a system \emph{would}
behave under a range of configurations, and it is offline. That difference is what makes
exhaustive construction of the transition graph, model checking, and the complex codomain
of Chapter~\ref{sec:quantum} viable here and not in an interpreter, and it is why the
endofunctor and the coalgebra are separated above. The endofunctor adjoins the map. The
coalgebra prices it.
\end{remark}

\begin{theorem}[Markovianity on configurations]
\label{wt:thm:markov}
Sampling $\gamma$ in the standard way --- draw a waiting time exponentially with parameter
$\propz(\mathfrak{c})$, select a redex with probability proportional to its contribution,
fire it --- yields a time-homogeneous continuous-time Markov chain on $\Cfg$, with
generator $Q(\mathfrak{c},\mathfrak{c}') = \gamma(\mathfrak{c})(\mathfrak{c}')$ off the
diagonal and $Q(\mathfrak{c},\mathfrak{c}) = -\sum_{\mathfrak{c}' \neq \mathfrak{c}}
Q(\mathfrak{c},\mathfrak{c}')$. It is \emph{not} in general a Markov chain on terms alone.
\end{theorem}

\begin{proof}[Proof sketch]
The propensity is a function of $\mathfrak{c}$ and of nothing else --- in particular not
of the time or the history --- and the successor is determined by $\mathfrak{c}$ together
with the selected redex. That is the definition of a time-homogeneous chain with the
stated generator. For the negative claim it suffices to exhibit two runs reaching the same
term with different maps, which Example~\ref{wt:ex:nonmarkov} does.
\end{proof}

\begin{example}[Dynamic weights are not a term-level phenomenon]
\label{wt:ex:nonmarkov}
Take rho with a single communication rule and keys $\varphi_a, \varphi_b, \varphi_{\bot}$,
where $\varphi_x$ says the redex is a communication on $x$. Let the updates act on
$\dec(\varphi_a)$ alone, by $u_a : w \mapsto w+1$ and $u_b : w \mapsto 2w$, which do not
commute. A term offering exactly one $a$-redex and one $b$-redex, both of which must fire,
reaches the same successor term by either order; starting from $\dec(\varphi_a) = 1$,
firing $a$ first arrives with $(1+1)\cdot 2 = 4$ and firing $b$ first with
$1\cdot 2 + 1 = 3$. The successor carries a single $a$-redex, so the two runs sit at the
same term with expected waiting times $1/4$ and $1/3$. No rate function of the term alone
reproduces both.
\end{example}

Theorem~\ref{wt:thm:markov} is the justification for
Definition~\ref{wt:def:config}, and the shape of the argument is worth being explicit
about, because the design could have gone otherwise. If weight maps could change but were
\emph{not} part of the state, the process on terms would be non-Markovian: the rate out of
$P$ would depend on how $P$ was reached. One would then be in the world of generalized
semi-Markov processes, with no exact simulation algorithm and no chain to model check.
Putting the map in the state restores exactness at the price of a larger state space, and
that price is real --- $\Cfg$ is $\terms \times \Semi^N$, a continuum unless the reachable
weight maps are finite in number. Remark~\ref{gw:rmk:state-is-term} will say what this
example looks like once there is no table.

\begin{remark}[Self-transitions]
\label{wt:rmk:selfloop}
Nothing forbids a rule whose firing returns a configuration $\equiv$-equal to the one it
fired in. Such a redex contributes to $\propz$ but to no off-diagonal entry, so the
diagonal above is $-\sum_{\mathfrak{c}' \neq \mathfrak{c}} Q$ and not $-\propz$. A
self-loop is a fictitious jump in the sense of uniformization: the two conventions induce
the same distribution of the state at every time and differ only in the number of recorded
events, hence in sojourn statistics. Neither is wrong; mixing them is.
\end{remark}

\subsection{What the restriction to names gives back}
\label{wt:sec:spim}

\begin{theorem}[Conservativity over summation-free SPiM]
\label{wt:thm:spim}
Restrict the construction to: the $\pi$-calculus fragment; keys that are channel-identity
predicates and nothing else; static maps, every update the identity; and every geometric
factor $1$. Then the propensity of Definition~\ref{wt:def:propensity} is exactly the
activity computed by the Stochastic Pi Machine for the summation-free
fragment~\cite{spim}, redex multiplicity included. The restriction is strict in each of
the four coordinates.
\end{theorem}

We do not reproduce the proof, which is a computation matching the two definitions of
activity term by term~\cite{StochSim}; the summation correction that the general SPiM
statement would need is empty here, since rho has no summation operator. The theorem is
worth having for what it says about the generality: everything the construction adds
beyond SPiM is visible as the removal of one of four restrictions, and each removal is
independently motivated. Formula-keys remove the first, plasticity the second, spatial
bias the third.

That completes the first construction. Everything above is correct, and everything
above is about to be shown to be a special case.

\section{Where the choice is}
\label{gw:sec:choice}

A rho program does not determine its own execution. Write
\begin{equation}
\label{gw:eq:race}
  x!(Q_1) \mid \textsf{for}(y \leftarrow x)\,P \mid x!(Q_2)
\end{equation}
and the operational semantics permits two reductions. It does not say which one happens,
and it does not say with what tendency. This is not an oversight. It is the defining
feature of a mobile process calculus, and it is what makes the calculus a specification
of interaction rather than of a machine.

An implementation must close the gap. RSpace --- the store that realizes rho in
f1r3node, to which this book returns when it turns from the calculus to the machine ---
closes it by presenting the
parallel fragment as a table from channels to bags of data and bags of continuations,
and letting the physical race between cores decide which produce is zipped with which
consume. That is a \emph{resolution mechanism}, and it is a choice, in the precise sense
Chapter~\ref{ch:chs} will make exact: the family of admissible pairings, indexed by
contended channels, has a section only when something supplies one.

The observation on which this half of the chapter rests is that the resolution mechanism
is a \emph{parameter}, and that other values of the parameter are available. Gillespie's
algorithm is a resolution mechanism: it supplies rates for the branches of a race and an
exponential clock for turning them into outcomes. The Born rule is a resolution mechanism:
it supplies amplitudes and a squared modulus. The weight map of the preceding sections is
a way of \emph{setting} the parameter from outside the term --- a table beside the
program. The question this half answers is what syntax a programmer needs in order to set
it from inside.

\subsection{Why contexts cannot see the parameter}

Before building the syntax it is worth knowing what kind of thing the parameter is, and
the cleanest way to say it is to say what cannot detect it.

\begin{principle}[Contexts create apertures; they do not read them]
\label{gw:prin:aperture}
A rho context can change which races exist --- by adding contention, by consuming a
contender, by installing a new receive. It cannot observe how any race is resolved.
Consequently no property of the form ``for all rho contexts $C[-]$, \ldots'' can be
sensitive to the resolution mechanism.
\end{principle}

The reason is almost a tautology once stated: a rho context is itself a rho term, whose
own races are resolved by the same unspecified mechanism. Composing silence with silence
yields silence. The consequence is that the correctness criterion the rest of this part
of the book leans on --- full abstraction against the class of contexts, as in
Chapter~\ref{ch:obsext} --- is not the right instrument here. It is not that there are
too few contexts. It is that the discriminating power one wants is not of the kind
contexts have.

Two things survive. The first is a different correctness frame: the object of study is
the map
\[
  \textsf{Resolvers} \times \textsf{Terms} \longrightarrow
  \textsf{Distributions over outcomes},
\]
in which the RSpace race, the Gillespie clock and the Born rule are three points of the
first coordinate, and the interesting statements are \emph{separations} between resolvers
at a fixed term rather than between terms at a fixed resolver. The second is a one-sided
obligation that every instantiation must meet.

\begin{definition}[Support adequacy]
\label{gw:def:support}
Let $\Reach(P)$ be the set of normal forms operationally reachable from $P$ under the
unweighted semantics, and let $\sem{P}$ be a weighted denotation of $P$ with values in
$\Semi$. The denotation is \emph{support-adequate} when
\[
  \supp\,\sem{P} \;\subseteq\; \Reach(P).
\]
\end{definition}

Support adequacy says the resolver invents nothing the operational semantics forbids.
The converse inclusion is \emph{not} demanded, and for a reason that will matter in
Chapter~\ref{sec:quantum}: a normal form that some derivation reaches, whose contributions
cancel, is exactly what a resolver with additive inverses can produce. So the natural
two-sided condition is not merely unproven; in the one instance where it fails, its
failure is the physics.

\begin{remark}[Non-compositionality is not a defect]
\label{gw:rmk:noncompositional}
The store denotation of rho is a monoid homomorphism for parallel composition, because it
denotes the store. A denotation that resolves races cannot be, because composing $P$ with
$Q$ creates races that neither had alone. The cut is the interaction term, and a
dependence across the cut that neither side carries is the semantic image of a race across
the cut. This is the same observation as Chapter~\ref{ch:apr}'s: an aperture is a
property of the cut and of neither agent flanking it. It is recorded here because the
reader trained on denotational semantics will otherwise read the failure of
compositionality below as a mistake.
\end{remark}

\section{The graded where-clause}
\label{gw:sec:clause}

\subsection{Resolution algebras}

\begin{definition}[Resolution algebra]
\label{gw:def:algebra}
A \emph{resolution algebra} is a commutative semiring
$\Semi = (V, \sumv, \tensorv, 0, 1)$: $\sumv$ associative and commutative with unit $0$,
$\tensorv$ associative and commutative with unit $1$, $\tensorv$ distributing over
$\sumv$, and $0$ annihilating.
\end{definition}

The two operations are forced by what the logical connectives must mean once they take
values. \emph{Conjunction is $\tensorv$}: two conditions holding together compose
multiplicatively, which is also what makes a joined receive
$p_1 \leftarrow c_1 \ \&\ p_2 \leftarrow c_2$ combine its clauses correctly.
\emph{Disjunction is $\sumv$}: two ways for a match to be enabled add. This is the same
semiring that Section~\ref{wt:sec:functor} took as a codomain; what has changed is that
it is now the value space of a \emph{logic}, and the logic's connectives are its
operations.

\begin{principle}[Interference is a logical phenomenon]
\label{gw:prin:or}
Since $\sem{\varphi \vee \psi} = \sem{\varphi} \sumv \sem{\psi}$, in a resolution
algebra with additive inverses a disjunction may evaluate to $0$ though neither
disjunct does. Destructive interference is a \emph{false disjunction with true
disjuncts}. This is the whole of the additional structure, stated in the logic rather
than imposed on it from outside.
\end{principle}

\begin{table}[t]
\centering
\small
\begin{tabular}{@{}llll@{}}
\toprule
\textbf{Algebra} & $\sumv,\tensorv$ & \textbf{Resolver} & \textbf{Regime}\\
\midrule
$\Bool = (\{0,1\},\vee,\wedge)$ & $\vee$, $\wedge$ & the RSpace race & rholang exactly \\
$\Rnn = (\mathbb{R}_{\geq 0},+,\times)$ & $+$, $\times$ & Gillespie & stochastic \\
$\CC = (\mathbb{C},+,\times)$ & $+$, $\times$ & Born rule & interference \\
Viterbi $([0,1],\max,\times)$ & $\max$, $\times$ & greedy & most-likely path \\
Tropical $(\mathbb{R}\cup\{\infty\},\min,+)$ & $\min$, $+$ & least cost & cost-optimal path \\
\bottomrule
\end{tabular}
\caption{Instances of the resolution algebra. One syntax; the resolver is the choice of
$\Semi$. The last two are not developed as clause instances in this book; they are
listed because the tropical instance is the cost reading that Chapter~\ref{sec:weighted}
develops in the linear fragment, where it supplies the least-action principle.}
\label{gw:tab:algebras}
\end{table}

Negation is the casualty and it is worth saying so plainly. On $[0,1]$ the map
$v \mapsto 1 - v$ serves; on $\CC$ there is no canonical complement and De Morgan fails
against $(+,\times)$. The design therefore keeps two sorts of formula, with the crisp
sort embedded in the graded one.

\subsection{Two sorts, one slot}

Let the \emph{crisp sort} be the condition language the book already has: ground
formulae supplied by a language definition, the spatial connectives generated from the
term formers, the context-labeled modality $\langle K \rangle$, and full Boolean
structure.
\begin{align*}
\beta &::= \textsf{true} \mid \textsf{false} \mid g \mid \neg\beta \mid
  \beta \wedge \beta \mid \beta \vee \beta \mid \langle K \rangle\beta \mid
  \textsf{out}(n)\beta \mid \textsf{in}(n)\beta \mid \beta \mid \beta \mid @\beta
\intertext{and let the \emph{graded sort} be}
\psi &::= \crisp{\beta} \mid v \mid \psi \tensorv \psi \mid \psi \sumv \psi
  \mid \langle K \rangle_{\Semi}\,\psi \qquad (v \in V).
\end{align*}
Here $\crisp{\beta}$ is the indicator embedding, sending a satisfied $\beta$ to $1$ and
an unsatisfied one to $0$, and $v$ ranges over scalar constants of the algebra. The
graded sort has no negation and no fixed-point operator; the crisp sort keeps both.

\begin{remark}[One slot, not two]
\label{gw:rmk:oneslot}
It is tempting to give the for-comprehension two clauses, a Boolean gate and a graded
weight. It is better to give it one graded clause and recover the gate as its
\emph{support}: a communication is enabled exactly when its clause evaluates to something
other than $0$. This keeps the language unforked --- the Boolean where-clauses rholang
already has are the $\Bool$-valued fragment under $\Bool \hookrightarrow \Semi$ --- and
it makes support adequacy (Definition~\ref{gw:def:support}) hold by construction rather
than by proof, since the classical interpreter sees nothing of the value except whether
it vanishes.
\end{remark}

\subsection{Syntax}

The for-comprehension already carries a \texttt{where} clause, whose condition may
mention the variables bound by the preceding receipts. The change is only to the sort of
that condition.

\begin{lstlisting}
for( p11 <- c11 & ... & pm1 <- cm1 ;
     ... ;
     p1n <- c1n & ... & pmn <- cmn
     where psi ) P
\end{lstlisting}

\noindent
with \texttt{psi} a graded formula over the algebra $\Semi$ declared for the shard or
for the channel's language definition. Within \texttt{psi} the bound names $p_{ij}$ are
in scope, so the value may depend on what was matched. That one sentence is the whole of
the difference between a clause and a key: a key classifies a redex by a formula fixed in
advance; a clause computes a value from the match in front of it.

\subsection{Per-match evaluation}
\label{gw:sec:permatch}

The decisive point concerns \emph{when} the clause is evaluated.

\begin{definition}[Candidates]
\label{gw:def:cand}
Let $r$ be a receipt site in $P$ --- a for-comprehension together with the channels its
receipts name. A \emph{candidate} at $r$ is an assignment $\sigma$ of a distinct datum
occurrence to each pattern of each group, such that every pattern matches its assigned
datum. Write $\Cand(P,r)$ for the set of candidates and
$\Out(r,\sigma) \in \terms/{\equiv}$ for the residual term the firing of $r$ under
$\sigma$ produces.
\end{definition}

Candidates are indexed by \emph{occurrences}, not by terms. This is the multiplicity
factor of Section~\ref{wt:sec:propensity} appearing in its native form: two structurally
congruent messages on a channel are two candidates, not one, because parallel
composition is a multiset and structural congruence does not delete duplicates. An
implementation that keys data by content alone cannot represent $\Cand$ faithfully.

\begin{definition}[Graded clause valuation]
\label{gw:def:val}
A graded clause $\psi$ at a receipt site $r$ denotes a function
\[
  \sem{\psi}_r : \Cand(P,r) \longrightarrow V ,
\]
defined by structural recursion, with $\crisp{\beta}$ evaluated against the term $P$
under the substitution $\sigma$ carries, scalars constant, $\tensorv$ and $\sumv$
pointwise, and $\langle K \rangle_{\Semi}$ as in Section~\ref{gw:sec:modality}.
\end{definition}

So a clause is a \emph{vector} indexed by candidates, and --- once composed with
$\Out$ --- a \emph{matrix} over candidates and outcomes. That is the whole mechanism,
and it is worth pausing on how little it is. The weight map needed a key set, a partition
discipline, an admissibility judgment, a classification function and an update function.
The clause needs a semiring and a scope rule.

\subsection{The graded modality is a transfer operator}
\label{gw:sec:modality}

Let $\Semi\langle \terms/{\equiv}\rangle$ be the free $\Semi$-semimodule on terms modulo
structural congruence. Define the one-step operator of a term by summing its candidates:
\begin{equation}
\label{gw:eq:onestep}
  T(P) \;=\; \sum_{r}\ \sum_{\sigma \in \Cand(P,r)}
    \sem{\psi_r}_r(\sigma)\;\cdot\;\Out(r,\sigma) ,
\end{equation}
the outer sum over receipt sites and the coefficient arithmetic in $\Semi$. Then the
graded modality is
\[
  \sem{\langle K \rangle_{\Semi}\,\psi}(P)
  \;=\; \sum_{r,\sigma} \sem{\psi_r}_r(\sigma) \tensorv \sem{\psi}(\Out(r,\sigma)),
\]
which is $T$ acting on the vector of values of $\psi$. The graded modality is the
transfer matrix of the reduction relation, and the principle that the meaning of a term
is the sum of all its possible futures is exactly the statement that $\sem{-}$ is a fixed
point of $\langle K\rangle_{\Semi}$. The denotational principle and the graded logic are
one object seen twice.

This is also where the chapter's own restriction~\ref{wt:req:str} is repaid. A key could
not carry a behavioral modality because a rate depending on what a redex does \emph{next}
cannot be recomputed locally. A clause \emph{may} carry $\langle K\rangle_{\Semi}$, and
the price is exactly the one named there: such a clause is not local, and an
implementation must either bound its modal depth or re-evaluate it globally. The
restriction has not gone away; it has become a choice the programmer makes in the clause,
with a cost the ledger can see.

\subsection{Denotation and conservativity}

\begin{definition}[Graded denotation]
\label{gw:def:denotation}
Fix a resolution algebra $\Semi$ and a budget $\beta \in \mathbb{N}$. Let
$M := \Semi\langle \terms/{\equiv}\rangle$ with basis vectors written $\lvert P \rangle$.
The $\beta$-budgeted denotation is
\[
  \sem{P}_\beta \;:=\; T_{\!*}^{\,\beta}\,\lvert P \rangle
  \;=\;\!\!\sum_{h \,:\, P \Rightarrow^{\leq\beta} Q}\!\! w(h)\,\lvert Q\rangle ,
\]
where $T_*$ extends $T$ linearly, $h$ ranges over causal histories of length at most
$\beta$ ending at a normal form or exhausting the budget, and $w(h)$ is the
$\tensorv$-product of the clause values along $h$. At a normal form $Q$ the operator acts
as the identity, $T\lvert Q\rangle = \lvert Q\rangle$, so that normal forms are absorbing.
\end{definition}

\begin{remark}[Why a budget, and only one]
\label{gw:rmk:budget}
The definition is corecursive and does not bottom out on non-terminating terms. For
terminating $P$ the recursion is well-founded on the reduction tree and $\beta$ may be
taken to be its depth. In general one needs either a convergence hypothesis --- a
spectral radius condition on $T$, since $\CC$ has no order and Knaster--Tarski is
unavailable --- or a budget. We take the budget, and note that it is the budget this book
already charges: the risk ledger's single $\omega$ for reflection
(Chapter~\ref{ch:tower}), and the token budget that keeps the mortal scientist's inquiry
terminating (\chSci). These should be one budget doing several jobs rather than several
budgets.
\end{remark}

One denotation admits three readings, differing only in what is done at the end.
\emph{Support}: $\supp(\sem{P}_\beta)$, which is what the classical interpreter sees and,
by Remark~\ref{gw:rmk:oneslot}, all it sees. \emph{Stochastic}: for $\Semi = \Rnn$, the
coefficients normalized, which is the jump-chain marginal of the Gillespie reading.
\emph{Born}: for $\Semi = \CC$, the squared moduli normalized, which is
Chapter~\ref{sec:quantum}.

\begin{theorem}[Conservativity]
\label{gw:thm:conservativity}
Let $\Semi = \Bool$ and let every clause be the embedding $\crisp{\beta}$ of a crisp
condition. Then $\supp(\sem{P}_\beta)$ is exactly the set of normal forms reachable from
$P$ in at most $\beta$ steps under the operational semantics with Boolean
\texttt{where} clauses.
\end{theorem}

\begin{proof}[Proof sketch]
In $\Bool$, $\tensorv = \wedge$ and $\sumv = \vee$, so $w(h) = 1$ iff every step of $h$
has a satisfied clause, and the coefficient of $\lvert Q\rangle$ is the disjunction over
histories reaching $Q$. Candidates coincide with the admissible zips of the store
presentation, and the absorbing convention does not add normal forms. Hence the
coefficient is $1$ iff some admissible derivation reaches $Q$.
\end{proof}

\begin{corollary}[Support adequacy]
\label{gw:cor:adequacy}
For any $\Semi$, $\supp(\sem{P}_\beta) \subseteq \Reach_\beta(P)$, since a nonzero
coefficient requires at least one history with all clause values nonzero, and such a
history is admissible under the Boolean image of the clauses. Equality fails exactly when
some reachable coefficient is a sum of nonzero terms summing to $0$.
\end{corollary}

This is the payoff of Remark~\ref{gw:rmk:oneslot}, and it should be read as a safety
theorem. Adding grading cannot make a program do something it could not do; it can only
make it fail to do something it could. Everything the ecology part builds on top of the
clause inherits this: a learner whose hypotheses are clauses can be wrong, but it cannot
be wrong in a way the unweighted calculus forbids.

\section{Weight maps are the linear graded clauses}
\label{gw:sec:dnf}

The first half of this chapter can now be placed.

\begin{theorem}[Weight maps are a normal form]
\label{gw:thm:dnf}
Every weight map $(\varphi_i, v_i)_{i\leq k}$ on a rule is denoted by the graded clause
\[
  \psi \;=\; \bigoplus_{i=1}^{k} \; v_i \tensorv \crisp{\varphi_i},
\]
and conversely a graded clause of this shape denotes a weight map if and only if the
$\varphi_i$ partition the refinement of the rule's left-hand side. Graded clauses not of
this shape --- those with nested $\tensorv$ over overlapping conditions, or with a graded
modality --- are strictly more expressive.
\end{theorem}

\begin{proof}[Proof sketch]
Forwards: on a candidate satisfying $\varphi_j$ and no other, all but the $j$th summand
vanish and the value is $v_j$, and exhaustiveness of the partition ensures some summand
survives. Backwards: if two $\varphi$'s overlap on a candidate, the clause returns
$v_i \sumv v_j$, which is a well-defined value but not a table lookup; the partition
condition is exactly the condition making the sum a selection. Strictness: with
overlapping $\varphi_i$ the clause is a mixture no weight map assigns, and
$\langle K\rangle_{\Semi}$ inspects successors, which no key can.
\end{proof}

The theorem is the hinge of the chapter and it is worth saying what it does to each piece
of the first half. The weight map is the \emph{linear} fragment of the clause language:
the clauses that are sums of scalars times indicators over a partition. The
classification function $\class_r$ is the statement that on such a clause exactly one
summand is nonzero. The propensity formula of Definition~\ref{wt:def:propensity} is
Equation~\eqref{gw:eq:onestep} with the clause in normal form and $\Out$ grouped by
outcome. And the convention Section~\ref{wt:sec:partition} declined --- ``sum the values
of all satisfied keys'' --- is not a convention: it is what $\sumv$ means, and the
partition discipline is what was needed to stop it from ever being invoked.

\begin{corollary}[The partition discipline is unnecessary]
\label{gw:cor:nopartition}
Single-valuedness of a graded clause is automatic, a formula having one value. The
requirements that key sets partition the left-hand-side refinement
(Definition~\ref{wt:def:partition}), that admissibility be checked
(Definition~\ref{wt:def:reqs}), and that a default class complete them
(Remark~\ref{wt:rmk:default}), all exist to secure single-valuedness for a lookup and
have no analogue here.
\end{corollary}

This dissolves the obstruction of Section~\ref{wt:sec:closure}. The Boolean demand ---
that the target logic specification supply negation so that a partition can be stated
--- arose because a partition needs complements. With no partition to define, the demand
does not arise. Negation is still absent from the graded sort, but it is absent as a
design choice about the value algebra rather than as an unmet requirement, and the seam
Observation~\ref{wt:obs:closure} pointed at is no longer there to point at. The reader
should notice what kind of repair this is. Nothing was proved that the weight map could
not prove; a presupposition was removed, and the problem it generated left with it. The
Prestige will have occasion to recall the shape.

\begin{proposition}[Plasticity is free]
\label{gw:prop:plastic}
A dynamic weight map, whose table is revised after each rewrite by an update function
$u : \Wmaps \times \textsf{Match} \times \textsf{Trace} \to \Wmaps$, is denoted by a
static graded clause whenever $u$'s dependence on the trace is expressible as a crisp
formula over the current term. No update function is needed, because a clause is
re-evaluated against the current term at every candidate.
\end{proposition}

\begin{proof}[Proof sketch]
The clause mentions the term through its crisp subformulae; its value at a candidate in
the term $P'$ reached after some history is by definition computed against $P'$. A table
revised to depend on features of $P'$ therefore agrees with the clause
$\bigoplus_i v_i \tensorv \crisp{\varphi_i}$ in which the $\varphi_i$ mention those
features directly. Dependence on trace history beyond what is visible in $P'$ is not
expressible, which is the hypothesis.
\end{proof}

\begin{remark}[The state is the term]
\label{gw:rmk:state-is-term}
Proposition~\ref{gw:prop:plastic} has a hypothesis, and Example~\ref{wt:ex:nonmarkov}
is a case where it fails: the two orders of firing leave the same term and different
tables, and no formula over the term can tell the orders apart. The weight map's answer
was to enlarge the state to $(P,\dec)$ and prove Markovianity there. The clause's answer
is the one rho always gives: \emph{put it in the term.} A rate that must remember something
the term does not show is a rate that depends on a datum nobody wrote down; write it down
--- a message on a channel the clause can read --- and the clause reads it. The
non-commuting updates of the example become two rewrites that each consume and re-emit a
counter, and the chain is a chain on terms again, because the term now carries what the
table carried. In a calculus where the state \emph{is} the term, ``the map is in the
state'' and ``the map is in the term'' are the same sentence, and the configuration of
Definition~\ref{wt:def:config} was a way of saying it before the syntax existed to say it
directly. Chapter~\ref{ch:spiking} will make this move for synaptic weights, and the
finiteness it needs is supplied by the same device that supplies it there: a quantized
value on a bounded range.
\end{remark}

\begin{remark}[What was retired, and what was kept]
\label{gw:rmk:retired}
The update function of Definition~\ref{wt:def:augmented} was the one component of the
weight map with no compilable form: an arbitrary closure over time and a trace summary,
for which a declarative sub-language would have had to be designed. The graded clause
\emph{is} that sub-language, and it was needed for other reasons. The geometric factor
$\gfac$ and the funding gate $\chi$ are kept exactly as they were: the first because it
weights the \emph{address} of a firing rather than the match, which a clause at the site
cannot see; the second because, as Proposition~\ref{wt:prop:gate} says, cost is not a
weighting and should not be dressed as one. The endofunctor $\Wgt_{\Semi}$ is kept as the
name for the linear fragment, which is what the physics part of the book instantiates.
\end{remark}

\section{Contention, and what is normalized}
\label{gw:sec:contention}

One more notion is needed before the clause can be handed to the chapters that use it,
and it is the notion the weight map never had to name because the partition hid it.
Equation~\eqref{gw:eq:onestep} sums over receipt sites as well as over candidates, and
that outer sum conflates two different things: candidates at different sites may be
\emph{alternatives}, or they may be \emph{concurrent}, and only the first should be
weighed against each other.

\begin{definition}[Conflict, components, matchings]
\label{gw:def:contention}
Two candidates \emph{conflict} when they share a datum or share a receipt. A
\emph{contention component} at $P$ is a connected component of the conflict relation. A
\emph{step} of a component is one of its \emph{maximal matchings}: a maximal set of
pairwise non-conflicting candidates.
\end{definition}

The terminology is that of event structures, where conflict names events that cannot both
occur~\cite{Winskel}. A contention component is the unit of choice: within it, the
alternatives are its maximal matchings, which is the RSpace zip --- the bag of data and
the bag of patterns, admissibly paired --- and across components nothing is chosen
because everything fires. Two independent races are two components; the race
equation~\eqref{gw:eq:race} is one component with two matchings; a receipt with $m$
patterns and $m$ data is one component with $m!$ matchings.

\begin{proposition}[Components factor]
\label{gw:prop:eventfact}
Let $C_1,\dots,C_k$ be the contention components at $P$. The maximal matchings of $P$
are exactly the elements of $\prod_i \mathrm{Match}(C_i)$: a step chooses one matching
per component, independently, and all of them fire.
\end{proposition}

\begin{proof}
Candidates in distinct components neither share a datum nor share a receipt, so a set of
candidates is pairwise non-conflicting iff its restriction to each component is;
maximality is likewise componentwise.
\end{proof}

Why this matters is that \emph{normalization} --- turning values into odds --- happens
over a component and nowhere else. In the Gillespie reading of Section~\ref{wt:sec:spim}
this was invisible, because exponential clocks race globally and the total propensity
does the normalizing for free. In the \emph{discrete} stochastic reading, where a
clause's values are relative odds rather than rates, the odds of a matching are its weight
divided by the sum over the matchings of its component; and in the complex reading, the
norm that must be preserved is taken over a component after summing amplitudes over
matchings with a common outcome. Both of those sentences will be said again, with their
consequences, in Chapters~\ref{sec:quantum} and~\ref{sec:weighted}. Here it is enough to
record two honest costs.

\begin{remark}[Two costs]
\label{gw:rmk:costs}
Normalization is \emph{nonlocal}: a clause can be normalized only in the light of what
else is contending, so a normalized value is a property of a configuration and not of a
clause. And the sum over maximal matchings of a component is permanent-shaped, hence
$\#P$-hard in general --- which is a fact about the physics it models and bites an
implementation, so what a compiler checks must be a syntactic sufficient condition such
as a bound on component size. Neither cost arises in the linear fragment, because a
partition is a component structure imposed in advance. That is a fair statement of what
the weight map bought and what it cost: tractability, at the price of forbidding every
race the modeler did not foresee.
\end{remark}

\section{Where this is used}

The construction is generic over the theory: it applies to anything one can write down as
a language presentation, so one obtains a stochastic machine for the ambient calculus, for
a user-defined domain language, or for cost-accounted rho, by writing the theory rather
than by writing a simulator.

Five places in this book take it up, and the division of labor between the two halves of
the chapter is worth stating. Chapter~\ref{ch:do}, next, builds the first thing a clause can do
that a table cannot --- an operator that intervenes on a race, under which Pearl's calculus of
causation embeds exactly. Chapter~\ref{ch:spiking} runs the clause at the smallest
scale where the whole apparatus is visible at once, and finds a spiking neural network
with plasticity in it --- synaptic efficacy as a graded clause, and the same rewrite that
performs inference performing the learning. \chSci\ takes the clause as the \emph{form of
a hypothesis}: a mortal scientist's conjecture about its world is a graded clause, and
what it learns is a value in the fiber over that clause. Chapter~\ref{sec:weighted} asks
which semiring, in the linear fragment, and finds that the choice is what separates
nondeterminism from cost from probability from amplitude. And Chapter~\ref{sec:quantum}
asks what the complex instance means, and finds that the clause takes the one route the
weight map could not --- amplitudes attached to candidates rather than to classes ---
and meets, in Definition~\ref{gw:def:contention}, exactly the structure that decides
whether the result is a quantum state.

% ===================================================================
